CAGRA fp32／int8、IVF-PQ 和 IVF-RaBitQ \cite{cagra,pq,ivfrabitq} 使用各自的网格（\cref{tab:3-grids}）；FP32 brute force 为穷举搜索。存储和重建精度可能不同。因此，主前沿衡量的是各自网格内达到质量目标时可实现的速度，并附单独的显存统计。

\par\medskip\noindent\begin{minipage}{\linewidth}
\centering\captionsetup{hypcap=false}
\captionof{table}[E1 主要调参维度]{E1 的主要调参维度。精确逐点参数与覆盖扩展保留在原始记录中。网格有限，且不同方法的网格不同。}
\label{tab:3-grids}
\footnotesize\setlength{\tabcolsep}{4pt}
\begin{tabularx}{\linewidth}{@{}>{\raggedright\arraybackslash}p{.22\linewidth}>{\raggedright\arraybackslash}X@{}}
\toprule
家族 & 记录的搜索维度 \\
\midrule
JHQ & Probes 与 $\alpha$；初始 probes 为 64/128/256，预算为工作负载预算的 $0.5/1/2$ 倍；覆盖扩展至 probe 1,024 和更大／更小预算 \\
CAGRA fp32/int8 & 图度数 32/64；搜索宽度 1/2/4/8；中间 top-$k$ 容量 32--512 \\
IVF-PQ & 八位子码；$d=768$ 时编码 192/384/768 bytes，更高维为 384/768/1,536；fp16/fp32 查表；probes 16--512 \\
IVF-RaBitQ & 每维 4/8/9 位、支持的搜索模式、probes 16--1,024；额外 4,096 列表 OpenAI 索引 \\
GPU 暴力搜索 & FP32 平方 L2 穷举，共同主机边界 \\
\bottomrule
\end{tabularx}
\end{minipage}\par\medskip

支配关系与实测前沿定义为
\[
 a\succ b\iff R_a\geq R_b,\ Q_a\geq Q_b,\
 (R_a,Q_a)\ne(R_b,Q_b).
\]
\begin{equation}
Q_{\rm meas}(r)=\max\{Q_i:R_i\geq r\}.
\label{eq:3-measured-frontier}
\end{equation}
没有满足要求的实测点时，结果未定义。对相邻前沿点，
\begin{equation}
\log Q_{\rm int}(r)=(1-\lambda)\log Q_a+\lambda\log Q_b,
\qquad \lambda=\frac{r-R_a}{R_b-R_a}.
\label{eq:3-interpolated-frontier}
\end{equation}
在 $5\times10^{-4}$ 内，解析器吸附到测点；否则只做插值，不外推。小批量比较中满足阈值的端点仍可入选。接近持平时同时报告两个估计器；显存未做匹配。

失败／无效输出被排除；仅有重复 ID 标记不会导致排除。隔离重跑替代受污染的构建；跳过的 Stella batch-1,000 运行不产生测点（\cref{sec:a-e1-protocol}）。Int8 仍纳入比较（\cref{fig:7-frontiers}）；ADC／海报图使用各自范围。
